\documentclass[11pt]{article}
\usepackage[margin=1in]{geometry}
\usepackage{amsmath,amssymb}
\usepackage{booktabs}
\usepackage{graphicx}
\usepackage{xcolor}
\usepackage{natbib}
\usepackage[hidelinks]{hyperref}

\newcommand{\TODO}[1]{\textcolor{red}{[#1]}}
\newcommand{\NUM}{\textcolor{gray}{\texttt{--.--}}}

\title{Do Elicited Epistemic States Predict Error?\\
A Predictive-Validity Test of Three-Component Uncertainty in Large Language Models}

\author{Maikel Leyva-V\'azquez\thanks{Universidad Bolivariana del Ecuador; Universidad Bernardo O'Higgins, Chile. ORCID 0000-0001-7911-5879.}
\ \and \TODO{coauthors}}

\date{Draft v0.1 --- \today}

\begin{document}
\maketitle

\begin{abstract}
Several formalisms propose representing a language model's epistemic state with more
than one number: truth, indeterminacy and falsity as separate components. Prior work
has shown that models will readily emit such triples, and that the components frequently
sum to more than one. What has not been established is whether the elicited triple carries
\emph{operationally useful} information. We state and test the question directly: holding
the prediction target fixed at answer correctness, does an elicited $(T,I,F)$ triple predict
error better than simpler signals obtained from the same model? We compare five signal
families --- elicited $(T,I,F)$, verbalized scalar confidence, self-verification $P(\text{True})$,
token-level sequence probability, and discrete semantic entropy --- under a single evaluation
protocol, on \TODO{D} short-form question-answering datasets and \TODO{M} models. Because
verbalized signals are known to be sensitive to how they are requested, every elicited arm is
run under \TODO{E} distinct elicitation protocols and we report the range across protocols
rather than a single point. We evaluate discrimination (AUROC, AUPRC), calibration
(Brier, ECE), and the decision-relevant quantity: selective risk under a
distribution-free conformal abstention rule at a fixed target risk level.
The decision rule is fixed in advance and stated in Section~\ref{sec:prereg}: the triple is
retained only if it beats verbalized scalar confidence by a margin whose confidence interval
excludes zero, under every elicitation protocol. \TODO{One-sentence result once data are in.}
All prompts, raw generations, and analysis code are released.
\end{abstract}

\section{Introduction}
\label{sec:intro}

A language model asked a question it cannot answer has, in principle, several distinguishable
predicaments. It may lack the information. The question may be vague. The question may be
about a future contingency. The proposition may be paradoxical. A single scalar confidence
score compresses all of these into one number, and a normalized output distribution forces
whatever mass is withheld from one answer onto the others.

This observation has motivated a family of proposals that represent the model's state with a
triple rather than a scalar: degrees of truth, indeterminacy and falsity treated as separate
and not required to sum to one \citep{leyva2026breaking}. Empirically it is established that
models will emit such triples when asked, that the components frequently sum to more than one,
and that the phenomenon replicates across vendors \citep{mason2026tensors}. It is also
established, by the same replication, that the scalar triple is \emph{insufficient} to separate
paradox from ignorance from contingency: distinct predicaments can produce identical triples.

None of this answers the question a practitioner would ask. Suppose we accept that the triple
is a richer description. Does it \emph{do} anything? Concretely: if we want to know which of a
model's answers will turn out to be wrong, does the triple tell us more than simply asking the
model how confident it is?

That is the question this paper tests. We fix the target --- $Y = 1$ if the model's answer is
incorrect --- and require every candidate signal to predict the same $Y$ on the same items with
the same evaluation protocol. This turns a representational claim into a measurable one.

\paragraph{Why the comparison must include verbalized scalar confidence.}
The decisive baseline is not semantic entropy. It is asking the model for one number.
If a three-component elicitation does not beat ``how confident are you, 0--100?'', then the
additional components are contributing vocabulary rather than information, and the
representational argument has no empirical support at the level of prediction. To our
knowledge this comparison has not been run. We treat it as the primary contrast and
every other signal as a reference point.

\paragraph{Why protocol variation is part of the design, not a robustness check.}
Verbalized confidence is strongly sensitive to elicitation format: recent work reports that
verbalized confidence tracks prompt format rather than prediction quality, with systematic
overconfidence and strong scale dependence \citep{protocolsens2026}. Any comparison between an
elicited triple and an elicited scalar that uses one prompt per arm is confounded by prompt
choice. We therefore run every elicited arm under $E$ protocols that vary wording, scale, and
response format, and report the full range. If between-protocol variation for a given signal
exceeds between-signal variation, that is the finding.

\paragraph{Contributions.}
\begin{enumerate}
  \item A preregistered predictive-validity test of elicited three-component epistemic states
        against four baselines under a single fixed target and evaluation protocol
        (Sections~\ref{sec:method}--\ref{sec:prereg}).
  \item An explicit decisive contrast against verbalized scalar confidence, and an ablation
        isolating the contribution of the indeterminacy channel $I$ over $T$ and $F$ alone.
  \item A protocol-sensitivity design in which each elicited signal is characterized by a range
        over $E$ elicitation protocols rather than a point estimate.
  \item A decision-relevant evaluation: conformal abstention at a fixed target risk, reporting
        achieved coverage per signal, rather than discrimination metrics alone.
  \item Public release of prompts, raw generations, gradings, and analysis code.
\end{enumerate}

\paragraph{What we are not claiming.} We do not test whether three-component representations
are philosophically preferable, nor whether they describe internal states. We test one thing:
whether the elicited numbers predict error. A negative result would not refute the
representational argument; it would establish that the argument cannot currently be supported
on predictive grounds, which is worth knowing and which we report either way.

\section{Background and related work}
\label{sec:related}

\paragraph{Uncertainty quantification for LLMs.}
The dominant sampling-based signal is semantic entropy, which samples multiple generations and
measures dispersion over meaning-equivalence classes rather than surface strings. Self-verification
approaches ask the model to judge its own candidate answer, yielding $P(\text{True})$.
Sequence-level token probability remains a strong and cheap baseline where logprobs are exposed.
Conformal prediction supplies distribution-free, finite-sample guarantees and is now routinely
combined with these scores to produce accept/abstain rules \citep{acse2026}.
Benchmarks in this area include long-form QA suites and unified libraries for UQ comparison.

\paragraph{Verbalized confidence and its fragility.}
A parallel line asks the model to state its confidence in words or numbers. The consistent
finding is that this channel is degenerate under many conditions: instruction-tuned models
concentrate reported confidence in a narrow high band regardless of correctness, and the
reported number is governed by response scale and prompt structure \citep{protocolsens2026}.
Probes on hidden states and single-pass logit entropy discriminate correctness better than the
verbal channel on the same items. This literature is the reason our design treats elicitation
protocol as a factor.

\paragraph{Abstention and knowledge boundaries.}
A growing benchmark family evaluates whether models decline to answer when they should
\citep{know2guess2026}. These benchmarks establish that helpfulness-tuned models suppress
refusal. They do not compare representations of the underlying epistemic state, which is the
gap this paper addresses.

\paragraph{Multi-component uncertainty formalisms.}
Representing belief with more than one number has several lineages: Dempster--Shafer belief
functions, subjective logic, imprecise probability, and four-valued Belnap--Dunn logics, each
with an established treatment of ignorance and conflict. Neutrosophic $(T,I,F)$ belongs to this
family and is the specific instantiation tested here \citep{leyva2026breaking}. We take no
position on the relations among these formalisms; we evaluate one of them empirically, and the
protocol we describe applies unchanged to the others.

\paragraph{Positioning.} The closest prior work established that unconstrained elicitation
produces $T+I+F>1$ at high rates and that this replicates across vendors
\citep{leyva2026breaking,mason2026tensors}. Neither establishes predictive validity. That is the
contribution here.

\section{What would count as evidence}
\label{sec:evidence}

Let $x$ be a question, $\hat{a}(x)$ the model's answer under greedy decoding, and
$Y(x) = \mathbb{1}[\hat{a}(x) \text{ is incorrect}]$ under a fixed grading function.
A \emph{signal} is any map $s: x \mapsto \mathbb{R}$ intended to be monotone in the probability
that $Y=1$. Every signal in this study predicts the same $Y$, on the same items, from the same
model, and is evaluated with the same metrics on the same split.

For a triple $(T,I,F)$ we define three scalarizations, fixed in advance:
\begin{align}
s_{\mathrm{TF}} &= F - T, \label{eq:stf}\\
s_{\mathrm{TIF}} &= F + \tfrac{1}{2} I - T, \label{eq:stif}\\
s_{\mathrm{TIF}}^{\mathrm{lr}} &= \text{fitted logistic combination of } (T,I,F) \text{ on the calibration split.} \label{eq:slr}
\end{align}
Equation~\eqref{eq:stf} ignores $I$ entirely; \eqref{eq:stif} charges indeterminacy at a fixed
rate; \eqref{eq:slr} lets the data choose the weights and is the most favourable treatment of
the triple. The contrast $s_{\mathrm{TIF}}^{\mathrm{lr}}$ versus $s_{\mathrm{TF}}$ is the
ablation that isolates whether $I$ carries information at all. Fitting \eqref{eq:slr} uses only
the calibration split; all reported metrics are computed on the held-out test split.

\section{Method}
\label{sec:method}

\subsection{Signals}
\begin{description}
  \item[A. Elicited triple.] One additional call per item requesting $T$, $I$, $F$ in $[0,1]$
        with no normalization constraint, parsed to a triple. Scalarized by
        \eqref{eq:stf}--\eqref{eq:slr}.
  \item[B. Verbalized scalar confidence.] One additional call requesting a single integer
        $0$--$100$. \emph{Primary baseline.}
  \item[C. Self-verification $P(\text{True})$.] One additional call presenting the question and
        the model's own answer and asking whether it is true; score is the model's stated
        probability, or the token probability of ``True'' where logprobs are available.
  \item[D. Sequence probability.] Mean token logprob of the greedy answer, where the provider
        exposes logprobs. Reported as available; missing for closed models without logprob support.
  \item[E. Discrete semantic entropy.] $K$ samples at temperature $1.0$, normalized and clustered
        into meaning-equivalence classes by answer normalization (short-form QA), then Shannon
        entropy over cluster frequencies. \TODO{If budget allows, add an NLI-based clustering variant
        as a robustness arm and report both.}
\end{description}

Arms A, B and C are elicited and therefore run under all $E$ protocols. Arms D and E are
sampling-based and have no elicitation protocol.

\subsection{Elicitation protocols}
$E = 3$ protocols per elicited arm, varying three things simultaneously so that they are
distinguishable formats rather than paraphrases:
P1, terse numeric, no explanation requested;
P2, explanation first then numbers (reasoning before the score);
P3, alternative scale --- verbal labels mapped to a fixed numeric grid, presented in reverse
order relative to P1. Full texts in Appendix~\ref{app:prompts}.

\subsection{Datasets}
$D = 3$ short-form QA datasets with gold answers and answer aliases, chosen so that grading is
mechanical and does not require a judge model: \TODO{name the three; suggested: TriviaQA
(rc.nocontext validation), SciQ, and a Natural Questions open-domain subset}. Short-form is a
deliberate restriction: it makes $Y$ well defined without an LLM grader, removing a confound
that would otherwise contaminate every arm equally but unpredictably.
$N = 1000$ items per dataset, sampled with a fixed seed, split $50/50$ into calibration and test.

\subsection{Models}
$M = 4$ models spanning at least three vendors and two size classes.
\TODO{Fix the list at run time and record exact version strings; do not report a model family
without its version.}

\subsection{Grading}
$Y$ is computed by normalized exact match against the gold answer set (lowercase, strip articles
and punctuation, alias expansion). Items where the model refuses to answer are recorded
separately and excluded from the discrimination analysis, with the refusal rate reported per
model and protocol; a signal cannot be credited for predicting the error of an answer that was
never given.

\subsection{Cost}
Per model and dataset: $N$ greedy answers, $N \times E \times 3$ elicitation calls (arms A, B, C),
$N \times K$ samples for arm E. With $N=1000$, $E=3$, $K=10$, $D=3$, $M=4$ this is
approximately \TODO{compute exactly from the run log} calls. Report the realized total and cost.

\section{Evaluation}
\label{sec:eval}

\paragraph{Discrimination.} AUROC and AUPRC for predicting $Y=1$.
\paragraph{Calibration.} Brier score and expected calibration error, after mapping each signal to
$[0,1]$ by isotonic regression fitted on the calibration split only.
\paragraph{Selective prediction.} Risk--coverage curve and its area (AURC), plus accuracy at
fixed coverage levels $\{0.9, 0.8, 0.5\}$.
\paragraph{Conformal abstention.} For each signal, calibrate a threshold on the calibration split
to achieve selective risk $\alpha \in \{0.05, 0.10\}$, then report \emph{achieved} risk and
\emph{achieved} coverage on the test split. Coverage at a fixed guaranteed risk is the number a
practitioner uses, and it is the headline comparison.
\paragraph{Uncertainty.} Paired bootstrap over items, $B = 10{,}000$ resamples, for every
pairwise difference between signals; report the difference and its $95\%$ interval, not two
separate intervals.
\paragraph{Protocol range.} For each elicited signal, report $\min$, $\max$ and range across the
$E$ protocols, alongside a variance decomposition of the metric into between-signal,
between-protocol, between-model and between-dataset components.

\section{Preregistered hypotheses and decision rules}
\label{sec:prereg}

Fixed before data collection. Any deviation is reported in Section~\ref{sec:limits}.

\begin{description}
  \item[H1 (primary).] $s_{\mathrm{TIF}}^{\mathrm{lr}}$ achieves higher AUROC than verbalized
        scalar confidence (arm B). \emph{Retain the triple only if} the paired bootstrap
        $95\%$ interval for the difference excludes zero \emph{under all three protocols} and in
        a majority of model $\times$ dataset cells.
  \item[H2 (ablation).] $s_{\mathrm{TIF}}^{\mathrm{lr}}$ achieves higher AUROC than
        $s_{\mathrm{TF}}$. If the interval includes zero, we conclude that the $I$ channel adds
        no predictive information, whatever H1 shows.
  \item[H3 (reference).] Semantic entropy (arm E) achieves higher AUROC than all elicited arms.
        We expect this and it is not a failure of the triple; it bounds what elicitation can buy.
  \item[H4 (protocol).] The between-protocol range of AUROC for a given elicited signal is
        smaller than the between-signal difference at fixed protocol. If H4 fails, the ranking of
        elicited signals is not identifiable from this design and we report it as such rather
        than reporting a winner.
  \item[H5 (decision).] At guaranteed selective risk $\alpha = 0.10$, the triple yields higher
        achieved coverage than arm B.
\end{description}

\paragraph{Reporting commitment.} All five hypotheses are reported with effect sizes and
intervals regardless of direction. If H1 and H2 both fail, the paper's conclusion is that
elicited three-component states do not carry predictive information beyond a single elicited
number, and the title and abstract will say so.

\section{Results}
\label{sec:results}

\TODO{All tables below are structural placeholders. No numbers are entered until
\texttt{run\_experiment.py} and \texttt{analyze.py} have been executed on the full design.}

\begin{table}[h]
\centering
\caption{Main discrimination results, pooled over datasets. AUROC for predicting incorrectness,
with paired-bootstrap $95\%$ intervals. Elicited arms report the range over $E=3$ protocols.}
\label{tab:main}
\begin{tabular}{llccc}
\toprule
Arm & Signal & AUROC (P1) & AUROC (P2) & AUROC (P3) \\
\midrule
A & $s_{\mathrm{TIF}}^{\mathrm{lr}}$ & \NUM & \NUM & \NUM \\
A & $s_{\mathrm{TIF}}$ (fixed weights) & \NUM & \NUM & \NUM \\
A & $s_{\mathrm{TF}}$ (no $I$) & \NUM & \NUM & \NUM \\
B & Verbalized scalar confidence & \NUM & \NUM & \NUM \\
C & $P(\text{True})$ & \NUM & \NUM & \NUM \\
D & Sequence logprob & \multicolumn{3}{c}{\NUM} \\
E & Discrete semantic entropy & \multicolumn{3}{c}{\NUM} \\
\bottomrule
\end{tabular}
\end{table}

\begin{table}[h]
\centering
\caption{Primary contrast (H1) and ablation (H2): paired differences in AUROC with $95\%$
bootstrap intervals, per protocol.}
\label{tab:contrast}
\begin{tabular}{lccc}
\toprule
Contrast & P1 & P2 & P3 \\
\midrule
$s_{\mathrm{TIF}}^{\mathrm{lr}} - $ verbalized scalar (H1) & \NUM & \NUM & \NUM \\
$s_{\mathrm{TIF}}^{\mathrm{lr}} - s_{\mathrm{TF}}$ (H2) & \NUM & \NUM & \NUM \\
\bottomrule
\end{tabular}
\end{table}

\begin{table}[h]
\centering
\caption{Decision-relevant comparison (H5): achieved coverage at guaranteed selective risk
$\alpha=0.10$ on the test split.}
\label{tab:conformal}
\begin{tabular}{lccc}
\toprule
Signal & Achieved risk & Achieved coverage & Coverage $95\%$ CI \\
\midrule
$s_{\mathrm{TIF}}^{\mathrm{lr}}$ & \NUM & \NUM & \NUM \\
Verbalized scalar & \NUM & \NUM & \NUM \\
$P(\text{True})$ & \NUM & \NUM & \NUM \\
Semantic entropy & \NUM & \NUM & \NUM \\
\bottomrule
\end{tabular}
\end{table}

\begin{table}[h]
\centering
\caption{Variance decomposition of AUROC (H4): share attributable to each design factor.}
\label{tab:variance}
\begin{tabular}{lc}
\toprule
Component & Share of variance \\
\midrule
Between signals & \NUM \\
Between elicitation protocols & \NUM \\
Between models & \NUM \\
Between datasets & \NUM \\
Residual & \NUM \\
\bottomrule
\end{tabular}
\end{table}

\section{Discussion}
\label{sec:discussion}

\TODO{Write the branch that the data select. Both branches are drafted below so that the
conclusion is not chosen after seeing the numbers.}

\paragraph{If H1 and H2 hold.} Elicited three-component states carry predictive information
beyond a single elicited number, and the indeterminacy channel is doing work. The claim then
available is narrow and defensible: for short-form QA, asking a model to separate its
uncertainty into components yields a usable error predictor that outperforms asking for one
number, at a cost of one additional call. This does not establish that the components correspond
to internal states, and we do not claim it.

\paragraph{If H1 holds but H2 fails.} The triple beats scalar confidence, but the gain comes from
$T$ and $F$ rather than from $I$. The honest reading is that asking for two opposing quantities
is a better elicitation format than asking for one, which is a prompting result, not a
representational one.

\paragraph{If H1 fails.} Elicited triples do not predict error better than a single elicited
number. The reported prevalence of $T+I+F>1$ then describes a property of unconstrained
elicitation rather than a usable signal, and any downstream argument that builds on the triple
as a measurement instrument lacks predictive support. We would state this in the abstract.

\paragraph{If H4 fails.} Protocol dominates signal. The practical recommendation would be that
comparisons between elicited uncertainty representations are not identifiable without protocol
randomization, and that published comparisons using a single prompt per arm should be read as
reporting a prompt effect.

\section{Limitations}
\label{sec:limits}

Short-form QA only; results may not transfer to long-form generation, where grading itself is
uncertain. Discrete semantic entropy over normalized strings is a weaker instrument than
NLI-based clustering. Sequence logprobs are unavailable for some providers, so arm D is not
complete across models. Elicitation is behavioural: nothing here licenses claims about internal
representations. $E=3$ protocols is a small sample from a large space; the protocol range is a
lower bound on protocol sensitivity, not an estimate of it. \TODO{Record every deviation from
Section~\ref{sec:prereg} here, including any that favour the hypothesis.}

\section{Conclusion}
\TODO{One paragraph, written after the analysis, stating which hypotheses survived.}

\section*{Reproducibility statement}
Prompts, raw generations, gradings, per-item scores, and analysis scripts are released at
\TODO{repository URL}. The analysis plan in Section~\ref{sec:prereg} was fixed and timestamped
before data collection; the timestamped file is included in the repository.

\appendix
\section{Elicitation prompts}
\label{app:prompts}
\TODO{Paste the exact three protocol texts from \texttt{prompts.py} verbatim, including system
messages and parsing regexes.}

\bibliographystyle{plainnat}
\bibliography{refs}

\end{document}
